% ID: MAT-GAN-RET-001
% Titolo: Retta passante per due punti
% Disciplina: Matematica
% Area: Geometria analitica
% Argomento: Retta
% Sottoargomento: Equazione della retta dati due punti
% Classe: Terza liceo scientifico
% Difficolta: 2
% Tipo: problema_numerico
% Risultato: 2x - 3y + 7 = 0; il punto C appartiene alla retta
% Tempo_stimato: 7 min
% Competenze: calcolo del coefficiente angolare, scrittura dell'equazione della retta, verifica di appartenenza
% Prerequisiti: piano cartesiano, coefficiente angolare, equazione della retta
% Tag: matematica, geometria_analitica, retta, coefficiente_angolare, appartenenza
% Fonte: esempio_originale
% Autore: OpenAI Codex
% Licenza: CC-BY-SA-4.0

\begin{esercizio}
Determina l'equazione della retta passante per i punti
\(A(-2,1)\) e \(B(4,5)\).

Verifica poi se il punto \(C(1,3)\) appartiene alla retta trovata.
\end{esercizio}

\begin{soluzione}
Calcoliamo il coefficiente angolare:
\[
m=\frac{y_B-y_A}{x_B-x_A}
=\frac{5-1}{4-(-2)}
=\frac{4}{6}
=\frac{2}{3}.
\]

Usiamo la forma punto-pendenza con il punto \(A(-2,1)\):
\[
y-1=\frac{2}{3}(x+2).
\]

Moltiplichiamo per \(3\):
\[
3y-3=2x+4.
\]

Portando tutto al primo membro:
\[
2x-3y+7=0.
\]

Per verificare se \(C(1,3)\) appartiene alla retta, sostituiamo le sue
coordinate:
\[
2\cdot 1-3\cdot 3+7=2-9+7=0.
\]

Poiche otteniamo \(0=0\), il punto \(C\) appartiene alla retta.
\end{soluzione}
